\section{Описание разработки проекта}
\label{sec:implementation}

\subsection{Архитектура программного комплекса}

Проект реализован на языке Haskell с использованием GHC \texttt{9.6.7} и системы сборки Stack \texttt{3.7.1}. Исходный код разделён на модули:

\begin{itemize}
    \item \texttt{Input.CLI} --- консольное меню, ввод пути к анализируемому файлу, выбор параметров анализа и запуск основного сценария.
    \item \texttt{Input.BDF} --- чтение BDF-файла и разбор его содержимого в структуру \texttt{BDFFont}.
    \item \texttt{Config.App} --- конфигурация приложения: путь отчёта по умолчанию, настройки логирования, пороги обнаружения аномалий и путь к эталонному шрифту.
    \item \texttt{Analyze.Metrics} --- расчёт метрик глифа и шрифта.
    \item \texttt{Diagnose.Anomaly} --- классификация аномалий по рассчитанным метрикам.
    \item \texttt{Diagnose.Replacement} --- поиск подходящей замены в эталонном шрифте.
    \item \texttt{Reference.Database} --- построение базы эталонных глифов из заданного в Config.App шрифта.
    \item \texttt{Report} --- сборка и форматирование итогового текстового отчёта.
    \item \texttt{Output} --- вывод отчёта и ошибок, запись отчёта в файл.
    \item \texttt{Logging} --- запись событий работы программы в журнал.
\end{itemize}

\subsection{Технические особенности реализации}

\subsubsection{Разбор BDF-файла}

Модуль \texttt{Input.BDF} читает файл как текст и выделяет блоки глифов между строками \texttt{STARTCHAR} и \texttt{ENDCHAR}. Для каждого глифа извлекаются имя символа, числовой код из строки \texttt{ENCODING} и строки пиксельной матрицы после маркера \texttt{BITMAP}. Если в файле отсутствует имя шрифта, код глифа или изображение глифа, функция возвращает ошибку типа \texttt{ParseError}.

\subsubsection{Метрики глифов}

Для каждого глифа рассчитываются четыре метрики:

\begin{itemize}
    \item Читаемость --- оценка формы глифа. Она учитывает баланс плотности, долю главной связной компоненты и покрытие минимальной прямоугольной области, содержащей закрашенные пиксели. Для поиска связных компонент используется алгоритм заливки области\footnote{пример реализации алгоритма заливки области на Haskell приведён в Rosetta Code.}.
    \item Пропорциональность --- отношение ширины пиксельной матрицы к высоте. Метрика показывает, насколько форма глифа вытянута по горизонтали или вертикали. Сама идея анализа размеров глифа опирается на принятые в шрифтовых форматах понятия растрового изображения глифа\footnote{\url{https://freetype.org/freetype2/docs/glyphs/glyphs-3.html}}.
    \item Плотность --- доля закрашенных пикселей в изображении глифа. Строки пиксельной матрицы берутся из BDF-представления, где изображение глифа задаётся шестнадцатеричными строками\footnote{\url{https://adobe-type-tools.github.io/font-tech-notes/pdfs/5005.BDF_Spec.pdf}; \url{https://font.tomchen.org/bdf_spec/examples/}.}.
    \item Различимость --- расстояние до ближайшего визуально похожего глифа. Сравнение основано на модифицированной мере Хэмминга для растровых представлений\footnote{\url{https://www.researchgate.net/publication/280561981_A_Modified_Hamming_Distance_Measure_for_Quick_Detection_of_Dissimilar_Binary_Images} --- Modified Hamming Distance Measure.}.
\end{itemize}

Пороговые значения выбраны эвристически для растровых шрифтов малого размера. Низкая плотность может означать почти пустой символ, слишком высокая плотность --- чрезмерно закрашенный символ, а малая различимость --- ситуацию, когда два разных глифа отличаются слишком малым количеством пикселей.

\subsubsection{Обнаружение аномалий}

Модуль \texttt{Diagnose.Anomaly} получает рассчитанные метрики и сравнивает их с порогами из \texttt{DetectionThresholds}. Глиф считается аномальным, если хотя бы один включённый параметр анализа вышел за допустимые границы:

\begin{itemize}
    \item читаемость меньше \texttt{0.4};
    \item плотность меньше \texttt{0.05};
    \item плотность больше \texttt{0.85};
    \item различимость меньше \texttt{0.016} с учётом допуска округления.
\end{itemize}

Значения \texttt{0.4}, \texttt{0.05} и \texttt{0.85} выбраны как эвристические границы для растровых шрифтов малого размера. Порог \texttt{0.016} для различимости был подобран вручную при формировании отчётов.

Пропорциональность выводится в отчёт как диагностическая метрика, но отдельный порог для неё в текущей конфигурации не задан.

\subsubsection{Поиск замен}

После обнаружения аномалий программа строит базу эталонных глифов из шрифта, заданного в конфигурации. По умолчанию используется шрифт Terminus. Для каждой аномалии модуль \texttt{Diagnose.Replacement} ищет глиф с таким же кодом символа. Если таких глифов нет, поиск выполняется по всей базе. Сходство считается как \texttt{1 - glyphDistance}. В отчёт попадает предложенная замена и значение сходства.

\subsubsection{Формирование отчёта}

Модуль \texttt{Report} собирает итоговый отчёт из заголовка, описания параметров анализа, таблицы метрик, списка аномалий и списка предложенных замен. Для удобства проверки рядом с числовыми значениями выводится пиксельное изображение каждого глифа. В разделе замен дополнительно выводятся два изображения: исходный глиф анализируемого шрифта и соответствующий глиф эталонного шрифта.

\subsubsection{Логирование}

Логирование реализовано в модуле \texttt{Logging}. События записываются в файл \texttt{coursework.log}. Каждая запись содержит время, уровень логирования и текст события. В журнал попадают запуск анализа, изменение параметров, ошибки чтения файлов и успешное сохранение отчёта.

\subsubsection{Тестирование свойств}

Для проверки ключевых свойств используется QuickCheck. Тесты проверяют, что плотность и читаемость находятся в диапазоне от 0 до 1, анализатор рассчитывает ровно одну запись метрик для каждого глифа, а нормальные значения метрик не приводят к появлению аномалий.

\newpage
